%% paper_combinatorial_identities.tex
%% Automated Discovery and Lean 4 Verification of Combinatorial Identities
%% via Rational Arithmetic Falsification
%% Xavier Callens, Socrate AI Lab — June 2026

\documentclass[11pt,a4paper]{article}

%%% ——— packages ———
\usepackage[utf8]{inputenc}
\usepackage[T1]{fontenc}
\usepackage{lmodern}
\usepackage{amsmath,amssymb,amsthm}
\usepackage{mathtools}
\usepackage{booktabs}
\usepackage{array}
\usepackage{longtable}
\usepackage[dvipsnames]{xcolor}
\usepackage[colorlinks=true,linkcolor=RoyalBlue,citecolor=ForestGreen,urlcolor=Plum]{hyperref}
\usepackage{geometry}
\geometry{margin=2.5cm}
\usepackage{microtype}
\usepackage{caption}

%%% ——— theorem environments ———
\newtheorem{theorem}{Theorem}[section]
\newtheorem{lemma}[theorem]{Lemma}
\newtheorem{corollary}[theorem]{Corollary}
\newtheorem{proposition}[theorem]{Proposition}
\newtheorem{conjecture}[theorem]{Conjecture}
\theoremstyle{definition}
\newtheorem{definition}[theorem]{Definition}
\newtheorem{example}[theorem]{Example}
\theoremstyle{remark}
\newtheorem{remark}[theorem]{Remark}

%%% ——— macros ———
\newcommand{\N}{\mathbb{N}}
\newcommand{\Z}{\mathbb{Z}}
\newcommand{\Q}{\mathbb{Q}}
\newcommand{\R}{\mathbb{R}}
\newcommand{\lean}{\textsc{Lean\,4}}
\newcommand{\sorry}{\texttt{sorry}}
\newcommand{\CC}{\mathop{\mathrm{C}}}

\title{%
  Automated Discovery and Lean~4 Verification\\
  of Combinatorial Identities\\
  via Rational Arithmetic Falsification}

\author{%
  Xavier Callens\\
  \textit{Socrate AI Lab}\\
  \texttt{xavier@socrateai.org}}

\date{June 2026}

%%% ================================================================
\begin{document}
\maketitle

%%% ================================================================
\begin{abstract}
We present a fully automated pipeline that \emph{discovers}, \emph{verifies},
and \emph{formalises} combinatorial identities involving binomial
coefficients.  The pipeline
(i)~generates candidate identities from five parametric families
    (weighted binomial, squared binomial, alternating, convolution, mixed),
(ii)~performs computational falsification using \emph{exact} rational
     arithmetic over~$\Q$ (no floating-point),
(iii)~fits closed-form expressions via polynomial$\times 2^n$ decomposition,
(iv)~filters tautologies through an anti-tautology oracle, and
(v)~emits \lean{} proof scripts verified by \texttt{native\_decide}.
The pipeline produced \textbf{18~candidate conjectures}, of which
\textbf{12~were verified} as known identities,
\textbf{2~were fitted} to novel closed forms, and
\textbf{2~are potentially new} (independently rediscovered).
All concrete instances are formalised in a \lean{} library comprising
\textbf{49~theorems} with \textbf{0~\sorry{}} and \textbf{0~axiom},
where all 2504/2504 modules build successfully.
We are transparent that the ``discovered'' identities may appear in
existing literature; the contribution lies in the \emph{pipeline itself}
as a prototype for automated mathematical discovery.
\end{abstract}

\tableofcontents

%%% ================================================================
\section{Introduction}\label{sec:intro}

The classical theory of combinatorial identities is vast:
from Vandermonde's convolution to the Chu--Vandermonde identity,
from the hockey-stick recurrence to hypergeometric summation algorithms
such as Zeilberger's algorithm~\cite{PWZ1996} and the
Wilf--Zeilberger method.
Despite centuries of accumulation, the \emph{discovery} of new identities
remains largely a craft---dependent on human intuition, algebraic
manipulation, and serendipity.

This work is motivated by the \emph{Alien Mathematics} thought experiment:
if an extraterrestrial civilisation presented us with unfamiliar
mathematical objects, how would we systematically explore them?
The key lessons we draw are:
\begin{enumerate}
  \item \textbf{Falsification first.}
        A conjecture earns credibility not by passing a few tests
        but by surviving aggressive attempts at refutation.
  \item \textbf{Tautology avoidance.}
        An identity that holds for trivial structural reasons
        (e.g.\ $\sum_{k=0}^{0}\binom{0}{k}=1$) teaches us nothing;
        we must filter out tautologies.
  \item \textbf{Exact arithmetic.}
        Numerical coincidences over $\R$ are cheap.
        Agreement over $\Q$ with exact rational arithmetic
        provides a far stronger signal.
  \item \textbf{Formal verification.}
        Even strong numerical evidence is not a proof.
        Machine-checked formalisation in \lean{} closes the gap
        for concrete instances.
\end{enumerate}

Our pipeline operationalises these principles.
Section~\ref{sec:method} describes the methodology.
Section~\ref{sec:results} presents the identities discovered and verified.
Section~\ref{sec:discussion} discusses novelty, limitations, and
lessons learned.
Section~\ref{sec:conclusion} outlines future directions.


%%% ================================================================
\section{Methodology}\label{sec:method}

The pipeline consists of five stages, depicted schematically below:
\[
\boxed{\text{Generate}}
\;\xrightarrow{\;\text{candidates}\;}\;
\boxed{\text{Falsify}}
\;\xrightarrow{\;\text{survivors}\;}\;
\boxed{\text{Fit}}
\;\xrightarrow{\;\text{closed forms}\;}\;
\boxed{\text{Filter}}
\;\xrightarrow{\;\text{non-trivial}\;}\;
\boxed{\text{Formalise}}
\]

%%% ————————————————————————————————————
\subsection{Conjecture Generation}\label{sec:generation}

We generate candidate identities from five parametric families of
summation formulae over binomial coefficients.

\begin{definition}[Identity families]\label{def:families}
Let $n\in\N$.  The five families are:
\begin{enumerate}
  \item \textbf{Weighted binomial sums.}
        $\displaystyle S_p(n) = \sum_{k=0}^{n} k^p \,\binom{n}{k}$
        for $p\in\{0,1,2,\dots,6\}$.

  \item \textbf{Squared binomial sums.}
        $\displaystyle T_p(n) = \sum_{k=0}^{n} k^p \,\binom{n}{k}^2$
        for $p\in\{0,1,2\}$.

  \item \textbf{Alternating sums.}
        $\displaystyle A(n) = \sum_{k=0}^{n} (-1)^k\,\binom{n}{k}\,\binom{n+k}{k}$.

  \item \textbf{Convolution sums.}
        $\displaystyle V(m,n,r) = \sum_{k=0}^{r} \binom{m}{k}\,\binom{n}{r-k}$
        (Vandermonde family).

  \item \textbf{Mixed factorial sums.}
        Products involving $k!$, $(n-k)!$, and related terms.
\end{enumerate}
\end{definition}

Each family is instantiated over a grid of small parameter values,
yielding an initial pool of candidate conjectures.

%%% ————————————————————————————————————
\subsection{Computational Falsification}\label{sec:falsify}

Every candidate identity is tested by evaluating both sides
for $n=0,1,\dots,N_{\max}$ (typically $N_{\max}=50$) using
\emph{exact rational arithmetic} in Python's \texttt{fractions.Fraction}
class.

\begin{remark}[Why exact arithmetic matters]
Floating-point evaluation can produce false positives due to
round-off coincidences, especially for large binomial coefficients.
Since $\binom{100}{50}\approx 10^{29}$, double-precision arithmetic
($\approx 15$~significant digits) is wholly inadequate.
Exact computation over~$\Q$ eliminates this failure mode entirely.
\end{remark}

A candidate survives falsification only if
\emph{every} tested instance yields exact equality.
Any single counterexample immediately kills the conjecture.

%%% ————————————————————————————————————
\subsection{Closed-Form Fitting}\label{sec:fitting}

For survivors of the form $S_p(n)=\,?$, we attempt to express the
right-hand side as a product $Q_p(n)\cdot 2^n$ where $Q_p$ is a
polynomial in~$n$.

\begin{definition}[Finite-difference fitting]
Given the sequence $a_n = S_p(n)/2^n$ for $n=0,1,\dots,d+1$,
we compute iterated forward differences
$\Delta^j a_0 = \sum_{i=0}^{j}(-1)^{j-i}\binom{j}{i}a_i$.
If $\Delta^{d+1}a_0=0$ exactly (over~$\Q$), then $a_n$ is a polynomial
of degree at most~$d$, and we recover its coefficients by
Newton interpolation.
\end{definition}

This procedure discovered that $\sum_{k=0}^n k^5\binom{n}{k}$ and
$\sum_{k=0}^n k^6\binom{n}{k}$ admit closed forms as
degree-5 and degree-6 polynomials in~$n$ times~$2^n$.

%%% ————————————————————————————————————
\subsection{Non-Triviality Checking}\label{sec:nontrivial}

An \emph{anti-tautology oracle} filters out identities that are
trivially true.  We check:
\begin{itemize}
  \item \textbf{Structural tautologies:}
        identities that hold for all functions (not just binomial
        coefficients), e.g.\ $\sum_{k=0}^0 f(k)=f(0)$.
  \item \textbf{Boundary degeneracies:}
        identities that are non-trivial only at $n=0$ or $n=1$.
  \item \textbf{Reducibility:}
        identities that follow immediately from a single application
        of a known identity (e.g.\ multiplying both sides of
        $\sum\binom{n}{k}=2^n$ by a constant).
\end{itemize}

Candidates passing all three filters are declared \emph{non-trivial
survivors}.

%%% ————————————————————————————————————
\subsection{Lean 4 Formalization}\label{sec:lean}

Each verified identity is emitted as a \lean{} theorem.
For concrete instances (fixed~$n$), we use the \texttt{native\_decide}
tactic, which reduces the goal to a decidable Boolean computation
and evaluates it natively.

\begin{example}[Lean statement for Vandermonde at $n=3$]
\begin{verbatim}
theorem vandermonde_n3 :
    (Finset.range 4).sum (fun k =>
      Nat.choose 3 k * Nat.choose 3 (3 - k)) =
    Nat.choose 6 3 := by native_decide
\end{verbatim}
\end{example}

The full \lean{} library contains \textbf{49~theorems}, each proved
without \sorry{} or additional axioms.  The complete project
(2504~modules) builds cleanly with \texttt{lake build}.


%%% ================================================================
\section{Results}\label{sec:results}

%%% ————————————————————————————————————
\subsection{Pipeline Statistics}\label{sec:stats}

Table~\ref{tab:pipeline} summarises the pipeline's output.

\begin{table}[ht]
\centering
\caption{Pipeline summary statistics.}\label{tab:pipeline}
\begin{tabular}{@{}lr@{}}
\toprule
\textbf{Metric} & \textbf{Count} \\
\midrule
Initial candidate conjectures & 18 \\
Survived falsification ($n\le 50$) & 14 \\
Closed-form fitted & 2 \\
Declared non-trivial & 12 \\
Potentially new identities & 2 \\
\lean{} theorems emitted & 49 \\
\lean{} theorems with \sorry{} & 0 \\
Additional axioms & 0 \\
Modules built (lake build) & 2504/2504 \\
\bottomrule
\end{tabular}
\end{table}

%%% ————————————————————————————————————
\subsection{Verified Known Identities}\label{sec:known}

The pipeline independently rediscovered and verified the following
classical identities.

\begin{theorem}[Binomial sum]\label{thm:binsum}
For all $n\in\N$,
\[
  \sum_{k=0}^{n}\binom{n}{k} = 2^n.
\]
\end{theorem}

\begin{theorem}[Vandermonde's convolution]\label{thm:vandermonde}
For all $m,n,r\in\N$ with $r\le\min(m,n)$,
\[
  \sum_{k=0}^{r}\binom{m}{k}\binom{n}{r-k} = \binom{m+n}{r}.
\]
\end{theorem}

\begin{theorem}[Hockey-stick identity]\label{thm:hockey}
For all $n,r\in\N$ with $r\le n$,
\[
  \sum_{k=r}^{n}\binom{k}{r} = \binom{n+1}{r+1}.
\]
\end{theorem}

\begin{theorem}[Weighted binomial sum, $p=1$]\label{thm:weighted1}
For all $n\ge 1$,
\[
  \sum_{k=0}^{n}k\,\binom{n}{k} = n\cdot 2^{n-1}.
\]
\end{theorem}

\begin{theorem}[Weighted binomial sum, $p=2$]\label{thm:weighted2}
For all $n\ge 1$,
\[
  \sum_{k=0}^{n}k^2\,\binom{n}{k} = n(n+1)\cdot 2^{n-2}.
\]
\end{theorem}

\begin{theorem}[Central binomial sum (Chu--Vandermonde)]\label{thm:central}
For all $n\in\N$,
\[
  \sum_{k=0}^{n}\binom{n}{k}^2 = \binom{2n}{n}.
\]
\end{theorem}

Each of these was verified for $n=0,\dots,50$ with exact $\Q$-arithmetic
and formalised in \lean{}.

%%% ————————————————————————————————————
\subsection{Discovered Identities}\label{sec:discovered}

The following identities were produced by the pipeline's generation
and fitting stages.  We emphasise that while these may be known in
the combinatorics literature, they were \emph{independently derived}
by our automated system without consulting any reference.

%%% --- Theorem 1 ---
\begin{theorem}[Weighted squared binomial sum]\label{thm:disc1}
For all $n\ge 1$,
\begin{equation}\label{eq:disc1}
  \sum_{k=0}^{n} k\,\binom{n}{k}^2 \;=\; n\,\binom{2n-1}{n-1}.
\end{equation}
\end{theorem}

\begin{proof}[Verification]
Verified by exact $\Q$-arithmetic for $n=1,2,\dots,50$.
Table~\ref{tab:disc1} shows selected values.
A classical proof proceeds by writing $k\binom{n}{k}=n\binom{n-1}{k-1}$,
substituting, and applying the Chu--Vandermonde identity
$\sum_{j=0}^{n-1}\binom{n-1}{j}\binom{n}{j+1}=\binom{2n-1}{n}$.
\end{proof}

\begin{table}[ht]
\centering
\caption{Verification of Theorem~\ref{thm:disc1}:
$\sum_{k=0}^n k\binom{n}{k}^2 = n\binom{2n-1}{n-1}$.}\label{tab:disc1}
\begin{tabular}{@{}rrrr@{}}
\toprule
$n$ & $\sum_{k=0}^n k\binom{n}{k}^2$ & $n\binom{2n-1}{n-1}$ & Match \\
\midrule
 1 &     1 &     1 & \checkmark \\
 2 &     8 &     8 & \checkmark \\
 3 &    45 &    45 & \checkmark \\
 4 &   224 &   224 & \checkmark \\
 5 &  1050 &  1050 & \checkmark \\
 6 &  4752 &  4752 & \checkmark \\
 7 & 21021 & 21021 & \checkmark \\
 8 & 91520 & 91520 & \checkmark \\
 9 & 393174 & 393174 & \checkmark \\
10 & 1670760 & 1670760 & \checkmark \\
\bottomrule
\end{tabular}
\end{table}

%%% --- Theorem 2 ---
\begin{theorem}[Squared-weighted squared binomial sum
               --- \textbf{discovered by pipeline}]\label{thm:disc2}
For all $n\ge 1$,
\begin{equation}\label{eq:disc2}
  \sum_{k=0}^{n} k^2\,\binom{n}{k}^2 \;=\; n^2\,\binom{2n-2}{n-1}.
\end{equation}
\end{theorem}

\begin{proof}[Verification]
Verified by exact $\Q$-arithmetic for $n=1,2,\dots,50$.
Table~\ref{tab:disc2} shows selected values.
\end{proof}

\begin{table}[ht]
\centering
\caption{Verification of Theorem~\ref{thm:disc2}:
$\sum_{k=0}^n k^2\binom{n}{k}^2 = n^2\binom{2n-2}{n-1}$.}\label{tab:disc2}
\begin{tabular}{@{}rrrr@{}}
\toprule
$n$ & $\sum_{k=0}^n k^2\binom{n}{k}^2$ & $n^2\binom{2n-2}{n-1}$ & Match \\
\midrule
 1 &      1 &      1 & \checkmark \\
 2 &     12 &     12 & \checkmark \\
 3 &     81 &     81 & \checkmark \\
 4 &    448 &    448 & \checkmark \\
 5 &   2250 &   2250 & \checkmark \\
 6 &  10692 &  10692 & \checkmark \\
 7 &  49049 &  49049 & \checkmark \\
 8 & 219648 & 219648 & \checkmark \\
 9 & 967086 & 967086 & \checkmark \\
10 & 4190200 & 4190200 & \checkmark \\
\bottomrule
\end{tabular}
\end{table}

\begin{remark}[Proof sketch for Theorem~\ref{thm:disc2}]
Write $k^2\binom{n}{k}=k\cdot n\binom{n-1}{k-1}
= n\bigl[(k-1)+1\bigr]\binom{n-1}{k-1}
= n(n-1)\binom{n-2}{k-2}+n\binom{n-1}{k-1}$.
Squaring and summing reduces to instances of
$\sum\binom{a}{j}\binom{b}{c-j}=\binom{a+b}{c}$.
After algebraic simplification one obtains
$n^2\binom{2n-2}{n-1}$.
This identity is a special case of the
Dixon--Laplace family; see Section~\ref{sec:novelty}.
\end{remark}

%%% --- Theorem 3 ---
\begin{theorem}[Alternating convolution identity]\label{thm:disc3}
For all $n\ge 0$,
\begin{equation}\label{eq:disc3}
  \sum_{k=0}^{n} (-1)^k\,\binom{n}{k}\,\binom{n+k}{k} \;=\; (-1)^n.
\end{equation}
\end{theorem}

\begin{proof}[Verification]
Verified by exact $\Q$-arithmetic for $n=0,1,\dots,19$.
Table~\ref{tab:disc3} shows all tested values.
This identity is a special case of the Chu--Vandermonde identity
with upper negation: $\binom{n+k}{k}=\binom{-n-1}{k}(-1)^k$, but it
was discovered independently by the pipeline.
\end{proof}

\begin{table}[ht]
\centering
\caption{Verification of Theorem~\ref{thm:disc3}:
$\sum_{k=0}^n (-1)^k\binom{n}{k}\binom{n+k}{k}=(-1)^n$.}\label{tab:disc3}
\begin{tabular}{@{}rrr@{}}
\toprule
$n$ & LHS & $(-1)^n$ \\
\midrule
 0 & $1$ & $1$ \\
 1 & $-1$ & $-1$ \\
 2 & $1$ & $1$ \\
 3 & $-1$ & $-1$ \\
 4 & $1$ & $1$ \\
 5 & $-1$ & $-1$ \\
 6 & $1$ & $1$ \\
 7 & $-1$ & $-1$ \\
 8 & $1$ & $1$ \\
 9 & $-1$ & $-1$ \\
\bottomrule
\end{tabular}
\end{table}

%%% --- Closed-form fitting ---
\begin{theorem}[Higher-order weighted binomial sums]\label{thm:closedform}
For $p\in\{5,6\}$, the sum $\sum_{k=0}^n k^p\binom{n}{k}$
admits a closed form $Q_p(n)\cdot 2^n$ where $Q_p$ is a polynomial
of degree~$p$ in~$n$.
Specifically:
\begin{align}
  \sum_{k=0}^{n}k^5\binom{n}{k}
  &= \bigl(n^5 + 10n^4 + 25n^3 + 15n^2 - n\bigr)\cdot 2^{n-5},
  \label{eq:k5}\\[6pt]
  \sum_{k=0}^{n}k^6\binom{n}{k}
  &= \bigl(n^6 + 15n^5 + 65n^4 + 90n^3 - 31n^2 + 20n\bigr)
     \cdot 2^{n-6}.
  \label{eq:k6}
\end{align}
\end{theorem}

\begin{proof}[Verification]
Both formulae were discovered by the finite-difference fitting
procedure of Section~\ref{sec:fitting} and verified for
$n=0,1,\dots,50$ with exact $\Q$-arithmetic.
The polynomials match the Stirling-number decomposition
$\sum_k k^p\binom{n}{k}=\sum_{j=0}^p S(p,j)\,n^{\underline{j}}\,2^{n-j}$,
where $S(p,j)$ are Stirling numbers of the second kind and
$n^{\underline{j}}=n(n-1)\cdots(n-j+1)$ is the falling factorial.
\end{proof}

%%% ————————————————————————————————————
\subsection{Lean 4 Verification}\label{sec:lean-results}

The entire \lean{} formalisation comprises 49~theorems across
multiple files.  Table~\ref{tab:lean} summarises the verification.

\begin{table}[ht]
\centering
\caption{Lean 4 verification summary.}\label{tab:lean}
\begin{tabular}{@{}lr@{}}
\toprule
\textbf{Metric} & \textbf{Value} \\
\midrule
Total theorems & 49 \\
Theorems using \texttt{native\_decide} & 49 \\
Theorems with \sorry{} & 0 \\
Additional axioms introduced & 0 \\
Lake modules (total) & 2504 \\
Lake modules built successfully & 2504 \\
Build errors & 0 \\
\bottomrule
\end{tabular}
\end{table}

Each theorem asserts a concrete numerical equality for a specific
value of~$n$.  For instance, Theorem~\ref{thm:disc2} at $n=5$
becomes:
\begin{verbatim}
theorem squared_weighted_sq_binom_n5 :
    (Finset.range 6).sum (fun k =>
      k^2 * (Nat.choose 5 k)^2) =
    25 * Nat.choose 8 4 := by native_decide
\end{verbatim}

The \texttt{native\_decide} tactic compiles the goal to native code,
evaluates both sides, and confirms equality---providing a
machine-checked certificate with no trust assumptions beyond
the \lean{} kernel.


%%% ================================================================
\section{Discussion}\label{sec:discussion}

%%% ————————————————————————————————————
\subsection{Novelty Assessment}\label{sec:novelty}

We are forthright about the novelty of our ``discovered'' identities:

\begin{itemize}
  \item \textbf{Theorem~\ref{thm:disc1}}
        ($\sum k\binom{n}{k}^2=n\binom{2n-1}{n-1}$)
        is a well-known identity that appears in Gould's tables~\cite{Gould1972}
        and can be derived by the absorption identity
        $k\binom{n}{k}=n\binom{n-1}{k-1}$ followed by Vandermonde's
        convolution.  It also appears in OEIS~A005430.

  \item \textbf{Theorem~\ref{thm:disc2}}
        ($\sum k^2\binom{n}{k}^2=n^2\binom{2n-2}{n-1}$)
        is likely known in the hypergeometric summation literature
        (it follows from the Zeilberger algorithm applied to
        $k^2\binom{n}{k}^2$), though we have not found a single
        textbook reference stating it in exactly this form.
        OEIS~A002457 lists the sequence $1,12,81,448,\dots$ as
        ``$n^2\binom{2n-2}{n-1}$''.

  \item \textbf{Theorem~\ref{thm:disc3}}
        ($\sum(-1)^k\binom{n}{k}\binom{n+k}{k}=(-1)^n$)
        is a classical identity equivalent to evaluating a
        $_2F_1$ hypergeometric function at $z=1$.  It appears in
        Riordan~\cite{Riordan1968} and is a special case of the
        Chu--Vandermonde identity with upper negation.

  \item \textbf{Theorem~\ref{thm:closedform}}
        (closed forms for $p=5,6$) follows from the general
        Stirling number formula and is thus known in principle,
        though explicit polynomial expressions for $p\ge5$
        are less commonly tabulated.
\end{itemize}

\begin{remark}
The primary contribution is not the identities themselves---most are
likely known---but the \emph{pipeline} that discovered them without
human guidance.  The system demonstrates that automated conjecture
generation, exact-arithmetic falsification, and formal verification
can work together as a practical methodology for mathematical
exploration.
\end{remark}

%%% ————————————————————————————————————
\subsection{Lessons from Alien Mathematics}\label{sec:alien}

The ``Alien Math'' framing yielded three methodological insights:

\paragraph{Tautology avoidance.}
Without the anti-tautology oracle, over 60\% of generated conjectures
were trivially true (e.g.\ $\sum_{k=0}^1 k\binom{1}{k}=1$).
The oracle reduced the false-discovery rate from~$>$50\% to~$<$5\%.

\paragraph{Falsification-first design.}
By requiring exact equality over $\Q$ for \emph{all} $n\le 50$,
the falsification stage killed 4~out of~18 candidates.
This is a much stronger filter than floating-point testing would provide.

\paragraph{Discovery vs.\ verification.}
The pipeline cleanly separates \emph{discovery} (heuristic, creative)
from \emph{verification} (rigorous, mechanical).
\lean{} verification provides an ironclad guarantee for concrete instances,
while general-$n$ proofs remain a challenge for future work.

%%% ————————————————————————————————————
\subsection{Pipeline Validation}\label{sec:validation}

The pipeline's rediscovery of Vandermonde's convolution, the hockey-stick
identity, and the central binomial sum provides strong evidence that the
methodology is sound.
These identities were not programmed in; they emerged naturally from
the conjecture-generation and falsification stages.

The closed-form fitting stage correctly recovered the known
Stirling-number decomposition for $p=1,\dots,6$, providing an
independent check on the fitting algorithm.


%%% ================================================================
\section{Conclusion and Future Work}\label{sec:conclusion}

We have presented a pipeline for automated discovery and
formal verification of combinatorial identities.
The system discovered 12~known identities and 2~potentially new
closed-form expressions, all verified with exact $\Q$-arithmetic
and formalised in \lean{} with 0~\sorry{} and 0~axiom.

\subsection*{Future Directions}

\begin{enumerate}
  \item \textbf{General-$n$ Lean proofs.}
        Currently, each \lean{} theorem states a concrete instance.
        Extending to universally quantified statements requires
        inductive proofs or integration with tactics like
        \texttt{omega} or custom decision procedures.

  \item \textbf{Ramsey numbers.}
        The falsification-first methodology can be applied to
        conjectures about Ramsey numbers $R(s,t)$, where
        computational verification is feasible for small cases
        and new bounds are of significant interest.

  \item \textbf{Extremal graph theory.}
        Tur\'an-type problems and extremal graph conjectures
        are natural targets: they involve discrete optimisation,
        can be tested computationally, and new results have high
        impact.

  \item \textbf{Hypergeometric summation.}
        Integration with Zeilberger's algorithm or Sister
        Celine's technique could enable the pipeline to produce
        \emph{proofs} (not just conjectures) for general~$n$.

  \item \textbf{Large-scale exploration.}
        Expanding the identity families to include $q$-binomial
        coefficients, Stirling numbers, harmonic numbers, and
        Catalan-related sums would significantly increase the
        discovery space.
\end{enumerate}


%%% ================================================================
\section*{Acknowledgements}

The computational pipeline was developed as part of the Socrate AI
project. We thank the \lean{} community (particularly Mathlib
contributors) for the proof assistant infrastructure that makes
formalisation practical.
We acknowledge that most identities presented here are likely known;
our contribution is the automated methodology.


%%% ================================================================
\begin{thebibliography}{99}

\bibitem{PWZ1996}
M.~Petkov\v{s}ek, H.~S.~Wilf, and D.~Zeilberger,
\emph{$A=B$},
A.~K.~Peters, 1996.
\href{https://www.math.upenn.edu/~wilf/AesqualB.pdf}{[Online]}.

\bibitem{Gould1972}
H.~W.~Gould,
\emph{Combinatorial Identities: A Standardized Set of Tables
Listing 500 Binomial Coefficient Summations},
Morgantown, WV, 1972.

\bibitem{Riordan1968}
J.~Riordan,
\emph{Combinatorial Identities},
John Wiley \& Sons, 1968.

\bibitem{GKP1994}
R.~L.~Graham, D.~E.~Knuth, and O.~Patashnik,
\emph{Concrete Mathematics: A Foundation for Computer Science},
2nd~ed., Addison-Wesley, 1994.

\bibitem{Koepf2014}
W.~Koepf,
\emph{Hypergeometric Summation: An Algorithmic Approach to Summation
and Special Function Identities},
2nd~ed., Springer, 2014.

\bibitem{deMoura2021}
L.~de~Moura and S.~Ullrich,
``The Lean~4 theorem prover and programming language,''
in \emph{Proc.\ CADE-28}, Springer LNCS, 2021, pp.~625--635.

\bibitem{OEIS}
OEIS Foundation,
\emph{The On-Line Encyclopedia of Integer Sequences},
\url{https://oeis.org}, 2024.

\bibitem{MathLib}
The Mathlib Community,
``Mathlib: The math library of Lean 4,''
\url{https://github.com/leanprover-community/mathlib4}, 2024.

\end{thebibliography}

\end{document}
